% Einheit 4 von 5 – Dreiecke konstruieren (bank/winkel-dreiecke/e4.jsonl, zusammenbau v0.1)
%% TODO Zweigzeile Teil 1: was hier gelernt wird (Satz in Schülersprache aus dem Einheitstitel) – Einheit 4
\einheitenkopf[e4]{Einheit 4 von 5 · Dreiecke konstruieren}
\zweigzeile{ab Kl. 6, je nach Buch bis Kl. 7 (am Gymnasium ab Kl. 5) · P10}
\setcounter{aufgabe}{19}

%% TODO Titel: Ich-kann-Satz fehlt in der Bank (Platzhalter: Kettenname) – Nr. 20
%% TODO Anweisung: ein Satz über der ersten Teilaufgabe fehlt in der Bank – Nr. 20
\begin{aufgabe}{Konstruieren}
\begin{teile}
\teil Gegeben sind a = 6 cm, b = 4 cm, $\gamma = 52^\circ$. Markiere die Stücke in einer Planfigur. Welcher Kongruenzsatz passt? \\ \kreuz{SSS} \\ \kreuz{SWS} \\ \kreuz{WSW} \\ \kreuz{SsW}
\end{teile}
\swa{Konstruiere das Dreieck ABC mit a = 6 cm, b = 5 cm und c = 7 cm. Beginne mit der Seite c auf dem Strahl.}{\winkelstrahl[9][A]}
\swa{Konstruiere das Dreieck ABC mit a = 5 cm, b = 4 cm und c = 6 cm. Beginne mit der Seite c auf dem Strahl.}{\winkelstrahl[9][A]}
\swa{Konstruiere das Dreieck ABC mit a = 8 cm, b = 6 cm und c = 7 cm. Beginne mit der Seite c auf dem Strahl.}{\winkelstrahl[9][A]}
\swa{Konstruiere das Dreieck ABC mit a = 4 cm, b = 7 cm und c = 5 cm. Beginne mit der Seite c auf dem Strahl.}{\winkelstrahl[9][A]}
\swa{Konstruiere das Dreieck ABC mit a = 6 cm, b = 7 cm und c = 8 cm. Beginne mit der Seite c auf dem Strahl.}{\winkelstrahl[9][A]}
\swfrage{Was bleibt gleich, was ändert sich?}
\swa{Konstruiere das Dreieck ABC mit b = 5 cm, c = 6 cm und $\alpha = 50^\circ$. Beginne mit der Seite c auf dem Strahl.}{\winkelstrahl[9][A]}
\swa{Konstruiere das Dreieck ABC mit c = 7 cm, $\alpha = 38^\circ$ und $\beta = 71^\circ$. Beginne mit der Seite c auf dem Strahl.}{\winkelstrahl[9][A]}
\swa{Konstruiere das Dreieck ABC mit b = 6 cm, c = 5 cm und $\beta = 64^\circ$. Beginne mit der Seite c auf dem Strahl.}{\winkelstrahl[9][A]}
\swa{Konstruiere ein gleichschenkliges Dreieck ABC mit der Basis c = 6 cm und den Schenkeln a = b = 5 cm.}{\winkelstrahl[9][A]}
\swa{Konstruiere ein rechtwinkliges Dreieck ABC mit dem rechten Winkel bei C und den Katheten a = 4 cm und b = 6 cm.}{\winkelstrahl[9][C]}
%% TODO Darstellung für schwach fehlt in der Bank (winkel-dreiecke-e4-k1-s7-v1; Abschnitt „Für schwache Schüler“)
\swz[3]{Schreibe eine Konstruktionsbeschreibung in drei bis vier Schritten für das Dreieck ABC mit c = 6 cm, $\alpha = 55^\circ$ und b = 4 cm.}{}
\begin{teile}
\teil Lässt sich aus den Seiten 4 cm, 5 cm und 10 cm ein Dreieck konstruieren? \\ \kreuz{ja} \\ \kreuz{nein}
\end{teile}
%% TODO Darstellung für schwach fehlt in der Bank (winkel-dreiecke-e4-k1-s9-v1; Abschnitt „Für schwache Schüler“)
\swz[3]{Dreieck ABC: AB = 5 cm, BC = 7 cm, CA = 6 cm. Dreieck PQR: PQ = 7 cm, QR = 6 cm, RP = 5 cm. Die Dreiecke sind kongruent. Welche Ecke von PQR gehört zu A, zu B und zu C?}{}
\begin{teile}
\teil Ein dreieckiges Beet ABC hat die Seite a = 12 m. Was gilt für die Summe der beiden anderen Seiten b und c? \hfill (P10 2020 OS) \\ \kreuz{$b + c > 12$ m} \\ \kreuz{$b + c = 12$ m} \\ \kreuz{$b + c < 12$ m}
\end{teile}
\end{aufgabe}

%% TODO Titel: Ich-kann-Satz fehlt in der Bank (Platzhalter: Kettenname) – Nr. 21
%% TODO Anweisung: ein Satz über der ersten Teilaufgabe fehlt in der Bank – Nr. 21
\begin{aufgabe}{Dreieck skizzieren und beschriften}
\swa{Beschrifte die Planfigur: Ecken A, B, C gegen den Uhrzeigersinn, dazu die Seiten a, b, c und die Winkel α, β, γ. Markiere farbig die gegebenen Stücke b = 5 cm, c = 6 cm und $\alpha = 38^\circ$.}{\dreieck{(0,0)}{(5,0)}{(1.5,3)}{}{}{}{}{}{}}
\end{aufgabe}

%% TODO Titel: Ich-kann-Satz fehlt in der Bank (Platzhalter: Kettenname) – Nr. 22
%% TODO Anweisung: ein Satz über der ersten Teilaufgabe fehlt in der Bank – Nr. 22
\begin{aufgabe}{Konstruieren – Fehler finden, Begründen, Anwendung}
\swz[3]{Sina sagt: „Mit $\alpha = 52^\circ$, $\beta = 64^\circ$ und $\gamma = 64^\circ$ ist das Dreieck eindeutig festgelegt – das ist der Kongruenzsatz WWW.“ Finde den Fehler.}{}
\swfrage{Warum reicht es nicht, von einem Dreieck nur die drei Winkel zu kennen? Begründe.}
\swz[3]{Ein dreieckiges Beet soll mit Brettern von 3,5 m, 2,0 m und 1,2 m Länge eingefasst werden. Geht das? Begründe.}{}
\end{aufgabe}

\uebersichtskasten{Beschriftung: Ecken A, B, C gegen den Uhrzeigersinn; Seite a liegt C gegenüber, b liegt B gegenüber, c liegt A gegenüber; Winkel α bei A, β bei B, γ bei C. \\ Ein Dreieck ist eindeutig, wenn drei passende Stücke bekannt sind: SSS (drei Seiten) · SWS (zwei Seiten und der Winkel dazwischen) · WSW (eine Seite und ihre beiden Winkel) · SsW (zwei Seiten und der Winkel gegenüber der längeren). Drei Winkel reichen nicht. \\ Ablauf: Planfigur mit den gegebenen Stücken → eine gegebene Seite zeichnen → Winkel antragen oder Kreisbögen schlagen → dritte Ecke, Seiten nachziehen, beschriften. \\ SSS mit a = 7 cm, b = 5 cm, c = 4 cm: c zeichnen, Kreisbogen um A mit 5 cm, Kreisbogen um B mit 7 cm, Schnittpunkt ist C. \\ Prüfen: Zwei Seiten zusammen müssen länger sein als die dritte (Dreiecksungleichung). \quad 3 cm, 4 cm, 8 cm: 3 + 4 < 8 – kein Dreieck.}
